\pdfoutput=1
\documentclass[11pt]{article}
\usepackage{mathfin-paper}

\title{The Fundamental Theorem of Asset Pricing,\\ Formalized in Lean~4}
\author{Raphael Coelho\thanks{Independent researcher. ORCID
  \href{https://orcid.org/0009-0001-6601-1023}{0009-0001-6601-1023}.
  Correspondence: \href{mailto:raphaelrrcoelho@gmail.com}{\texttt{raphaelrrcoelho@gmail.com}}.
  Artifact: \url{https://github.com/formal-applied-math/formal-mathfin}, release \texttt{v1.5.0}.}}
\date{October 2026}

\hypersetup{pdftitle={The Fundamental Theorem of Asset Pricing, Formalized in Lean 4},
  pdfauthor={Raphael Coelho}}

\begin{document}
\maketitle

\begin{abstract}
\noindent
The Fundamental Theorem of Asset Pricing states that a market is free of arbitrage exactly when it
admits an equivalent martingale measure. We formalize it in Lean~4 over Mathlib in three settings: a
finite-state market with one risky asset over a finite horizon (Harrison--Pliska), a one-period market with one risky asset
on an arbitrary probability space, and a one-period market with finitely many assets on an arbitrary
probability space. The finite case is the geometry of a separating hyperplane; the scalar one-period
case is an elementary change of density. In the $d$-asset case the proof builds the equivalent martingale
measure from a minimiser of the potential $\theta\mapsto\E[\log(1+e^{\langle\theta,Y\rangle})]$: absence of
arbitrage makes the potential coercive off the redundant directions, its first-order condition is the
martingale property, and for integrable returns the logistic weight of the minimiser, normalised, is the
density; a bounded tempering removes the integrability assumption, so the return need only be measurable. The
construction uses no Hahn--Banach theorem, no $L^0$-closedness argument, no measurable selection and no
non-redundancy hypothesis. The separation lemma behind the finite
case has a companion, built from the same two helper lemmas, that yields the Artzner--Delbaen--Eber--Heath
representation of coherent risk measures on a finite state space. The boundary is explicit. The general multi-period
theorem of Dalang, Morton and Willinger is formalized only in its forward direction, and so is the
continuous-time theorem, for simple strategies.
Every theorem is \texttt{sorry}-free, the axioms of each result discussed are pinned to Lean's three
standard axioms by a build-enforced gate, and the development is reproducible from a pinned
toolchain.
\end{abstract}

%======================================================================
\section{Introduction}

A market offers an \emph{arbitrage} if some admissible strategy turns no initial wealth into a payoff
that is non-negative almost surely and strictly positive with positive probability. The Fundamental
Theorem of Asset Pricing identifies the absence of such opportunities with a single analytic object: a
probability measure $Q$, equivalent to the physical measure $P$, under which discounted asset prices
are martingales. No arbitrage holds if and only if such an \emph{equivalent martingale measure}
exists. The equivalence was made precise by Harrison and Kreps~\cite{HarrisonKreps1979} and Harrison
and Pliska~\cite{HarrisonPliska1981}, established for general discrete-time models by Dalang, Morton
and Willinger~\cite{DMW1990}, and carried to continuous time in the no-free-lunch form of Delbaen and
Schachermayer~\cite{DelbaenSchachermayer1994}. It is the result on which arbitrage-free pricing rests:
a price is the $Q$-expectation of a discounted payoff.

The theorem is also a natural test for formal mathematics, because its content is analytic rather than
computational. In the finite case it hides a separating-hyperplane argument; in general it asserts the
existence of a change of measure that is usually obtained non-constructively, by separating a cone in
$L^1$ with a Hahn--Banach functional after proving the cone closed in $L^0$. Fixing exactly which
hypotheses each direction needs, and checking the measure-theoretic details against a kernel, is a task
at which formalization is decisive.

Formalized probability has grown substantially, and with it the measure theory on which pricing
depends. Echenim, Guiol and Peltier verified Cox--Ross--Rubinstein pricing in
Isabelle/HOL~\cite{EchenimGuiolPeltier2020}, Keskin formalized discrete- and continuous-time
martingales~\cite{Keskin2023}, Avigad, H\"olzl and Serafin the central limit
theorem~\cite{AvigadHolzlSerafin2017}, and Degenne and collaborators constructed Brownian motion in
Lean~\cite{DegenneEtAl2025,BrownianMotionRepo}. The nearest precedent for the present paper is
Nagy's Lean~4 derivation from It\^o's lemma to Black--Scholes~\cite{Nagy2026}, which contains the
theorem for a one-period market with one asset and two states. A search in October 2026 of the Lean
ecosystem, the Isabelle Archive of Formal Proofs and the Coq/Rocq ecosystem found no machine-checked
proof beyond that case: none on an arbitrary probability space, none with several assets on an
arbitrary space, and none over several periods.

\paragraph{What this paper does}
We formalize the theorem in three settings of increasing generality and, in the most general of them,
give the equivalent martingale measure by a formula.
\begin{enumerate}[label=(\arabic*)]
\item \emph{Finite states, finitely many periods} (\cref{sec:finite}). For one risky asset, the backward
  direction is the geometry of a hyperplane separating the attainable gains from the standard simplex
  (\lean{ftap_discrete}); in one period the same separation handles finitely many assets.
\item \emph{One period, one asset, arbitrary probability space} (\cref{sec:scalar}). Integrability is
  no longer free, and the measure is built by an elementary change of density
  (\lean{ftap_one_period}).
\item \emph{One period, $d$ assets, arbitrary probability space} (\cref{sec:vector}). For integrable
  returns the proof takes as density of the equivalent martingale measure the normalised logistic weight of
  a minimiser of the softplus potential $\E[\log(1+e^{\langle\theta,Y\rangle})]$; a bounded tempering removes
  integrability.
  Redundant assets are absorbed rather than excluded, and the return need only be measurable
  (\lean{ftap_one_period_vector}).
\item \emph{One separation for pricing and risk} (\cref{sec:separation}). The separation lemma of the
  finite case and the one behind the representation of coherent risk measures on a finite state space
  (\lean{coherentRisk_isLUB}) are two faces of one finite-dimensional separation, proved side by side
  from shared helper lemmas.
\item \emph{A precise boundary} (\cref{sec:boundary}). The general multi-period theorem and the
  continuous-time theorem are formalized only in their forward directions; we say exactly what is
  proved.
\end{enumerate}

This paper is part of a Lean~4~\cite{deMouraUllrich2021} library of formalized mathematical finance
built on Mathlib~\cite{Mathlib2020}; the library is described in a companion
overview~\cite{CoelhoLibrary} and its It\^o calculus in a second companion~\cite{CoelhoIto}. The
results developed here live entirely within Mathlib's measure theory and convex analysis, on its
\lean{withDensity} change of measure and its finite-dimensional separation theorems; they stand
alongside, but do not depend on, the library's Brownian-motion layer.

\paragraph{Changes from the first version}
This version corrects three statements of the first. The two-state case was formalized
earlier~\cite{Nagy2026}, so the claim of priority is narrowed. The $d$-asset density is the softplus
analogue of the Esscher measure, not the Esscher measure itself (\cref{sec:why-softplus}). And only one
direction of the correspondence between no arbitrage and coercivity is proved, so the text no longer
calls them equivalent. The version adds the separation shared with coherent risk measures
(\cref{sec:separation}), the finite-state $d$-asset theorem and the forward directions that border the
development (\cref{sec:finite,sec:boundary}), and updates the artifact to the current release
(\cref{sec:artifact}).

%======================================================================
\section{The market and the two easy facts}\label{sec:market}

In each setting a position in the traded assets produces a discounted gain, and an arbitrage is a
position whose gain is non-negative almost surely without being almost surely zero. The settings differ
only in what a position is and how its gain is formed.
\begin{itemize}
\item In the \emph{finite-state} model a filtration $\F=(\F_t)_{t\le T}$ on a finite probability space
  carries an adapted discounted price $S_0,\dots,S_T$ of one risky asset. A strategy $\varphi$ is
  predictable, with $\varphi_{t+1}$ measurable for $\F_t$, and its discounted gain is the martingale
  transform $(\varphi\cdot S)_T=\sum_{t<T}\varphi_{t+1}(S_{t+1}-S_t)$.
\item In the \emph{scalar one-period} model a discounted excess return $Y\colon\Omega\to\R$ is held in
  position $\theta\in\R$, with gain $\theta Y$.
\item In the \emph{$d$-asset one-period} model the return $Y\colon\Omega\to F$ takes values in a
  finite-dimensional real inner-product space $F$ (the $d$-asset market is $F=\R^d$), a constant
  position $\theta\in F$ is held, and the gain is $\langle\theta,Y\rangle$. (In Mathlib the $d$-asset
  space is \lean{EuclideanSpace ℝ (Fin d)}, not the product $\R^d$ with its sup norm.)
\end{itemize}
Write $G$ for the discounted gain of a position. \emph{No arbitrage} says that a gain which is
non-negative is null:
\[
  G\ge 0 \quad P\text{-a.s.} \qquad\Longrightarrow\qquad G=0\quad P\text{-a.s.},
\]
which is the textbook ``no $G\ge0$ with $P(G>0)>0$'' in contrapositive form. The condition is stated
under $P$ alone and names no second measure. An \emph{equivalent martingale measure} is a probability
measure $Q$ with $Q\ll P$ and $P\ll Q$ under which discounted prices are fair: in the one-period
settings $Y$ is $Q$-integrable with $\E_Q[Y]=0$ (vector-valued in the $d$-asset case); in the
finite-state setting $S_t=\E_Q[S_{t+1}\mid\F_t]$ for every $t<T$. Recording equivalence as the pair
of absolute continuities is what lets a $Q$-almost-sure statement transfer back to $P$.
\Cref{lst:defs} shows the $d$-asset definitions; the other two settings differ only in the gain.

\begin{leancode}[caption={No arbitrage and the equivalent martingale measure, $d$-asset one-period market (\texttt{Foundations/FTAPOnePeriodVector}).},label=lst:defs]
variable {Ω : Type*} {mΩ : MeasurableSpace Ω}
  (P : Measure Ω) [IsProbabilityMeasure P]
  {F : Type*} [NormedAddCommGroup F] [InnerProductSpace ℝ F]
  [FiniteDimensional ℝ F] [MeasurableSpace F] [BorelSpace F] (Y : Ω → F)

def NoArbitrage : Prop :=
  ∀ θ : F, 0 ≤ᵐ[P] (fun ω ↦ inner ℝ θ (Y ω)) →
    (fun ω ↦ inner ℝ θ (Y ω)) =ᵐ[P] 0

structure IsEMM (Q : Measure Ω) : Prop where
  prob : IsProbabilityMeasure Q
  absP : Q ≪ P
  Pabs : P ≪ Q
  int  : Integrable Y Q
  fair : ∫ ω, Y ω ∂Q = 0
\end{leancode}

One direction of the theorem is uniform and elementary. If an equivalent martingale measure $Q$
exists, the expected discounted gain of every position vanishes under $Q$: in the one-period settings
$\E_Q[\theta Y]=\theta\,\E_Q[Y]=0$ and $\E_Q[\langle\theta,Y\rangle]=\langle\theta,\E_Q[Y]\rangle=0$, and
in the finite-state setting $\E_Q[(\varphi\cdot S)_T]=0$ because each term $\varphi_{t+1}(S_{t+1}-S_t)$ has
zero $\F_t$-conditional mean. A gain
with $G\ge0$ and $\E_Q[G]=0$ is zero $Q$-almost surely, hence $P$-almost surely by equivalence, so the
position is not an arbitrage. This is \lean{noArbitrage_of_isEMM} in each development. The content of
the theorem is the converse, which the next three sections supply in rising generality.

%======================================================================
\section{Finite markets: the geometric theorem}\label{sec:finite}

In a finite-state market the equivalent martingale measure is, almost literally, a separating
hyperplane. Fix a finite probability space $(\Omega,P)$ with full support, a filtration $\F$, an adapted
scalar discounted price $S$ and a horizon $T$. The discounted gains attainable by predictable strategies
form a linear subspace of $\R^\Omega$, the span of the one-period gains
$\ind_A(S_{t+1}-S_t)$ with $t<T$ and $A\in\F_t$; call it the \emph{gains subspace}. No arbitrage says
exactly that this subspace meets the non-negative cone only at the origin, equivalently that it is
disjoint from the standard simplex of probability vectors on $\Omega$.

The simplex is compact and convex and the gains subspace is closed, being finite-dimensional, so the
geometric Hahn--Banach theorem separates them strictly. The library takes this one level higher, in a
form that serves more than pricing. Its root (\cref{lst:cone}) is stated for a closed convex
\emph{cone} $C\subseteq\R^\iota$ disjoint from the simplex, and produces a strictly positive functional
$q$ that is non-positive on $C$. A subspace is a cone closed under negation, so applied to the gains
subspace the one-sided bound becomes two-sided: $q$ annihilates every attainable gain
(\lean{exists_pos_dual_of_disjoint_stdSimplex}). Read as a measure, $q$ makes every martingale transform
have zero expectation.

\begin{leancode}[caption={The cone-separation root (\texttt{Foundations/ConvexDuality}).},label=lst:cone]
theorem exists_pos_separating_of_cone_disjoint_simplex
    {ι : Type*} [Fintype ι]
    (C : Set (ι → ℝ)) (hCconv : Convex ℝ C) (hCclosed : IsClosed C)
    (hCcone : ∀ ⦃x⦄, x ∈ C → ∀ ⦃c : ℝ⦄, 0 ≤ c → c • x ∈ C)
    (hdisj : Disjoint (stdSimplex ℝ ι) C) :
    ∃ q : ι → ℝ, (∀ i, 0 < q i) ∧ ∀ v ∈ C, ∑ i, q i * v i ≤ 0
\end{leancode}

The cone form is not needed for the market, whose gains form a subspace; it is stated for the sake
of \cref{sec:separation}, where a genuine cone appears. The market reaches the theorem in two steps. Under no arbitrage the gains subspace misses the simplex
(\lean{gains_disjoint_stdSimplex}): a non-negative gain is null, and on a full-support space a null
vector has no positive coordinate, while every point of the simplex has one. The strictly positive
annihilator, normalised to a probability vector, is a measure $Q\sim P$; it annihilates each generator
$\ind_A(S_{t+1}-S_t)$, which is the defining property of the conditional expectation, so $S$ is a
$Q$-martingale up to $T$. With the forward direction of \cref{sec:market} this is the biconditional
(\cref{lst:discrete}).

\begin{leancode}[caption={The finite-state theorem (\texttt{Foundations/FTAPDiscrete}). $\Omega$ is a \texttt{Fintype} with \texttt{MeasurableSingletonClass}, $\mathcal{F}$ a filtration indexed by $\mathbb{N}$, $S:\mathbb{N}\to\Omega\to\mathbb{R}$.},label=lst:discrete]
theorem ftap_discrete (hS : StronglyAdapted 𝓕 S) (hP : ∀ ω, 0 < P {ω}) :
    NoArbitrage 𝓕 P S T ↔ ∃ Q, IsEMM 𝓕 P S T Q
\end{leancode}

This is the theorem of Harrison and Pliska, the finite case of Dalang, Morton and Willinger. The
statement is slightly stronger than the textbook one: $P$ need only be a measure of full support on the
finite space, not a probability measure, and strategies are unbounded. The same
kernel gives the one-period theorem for finitely many assets on a finite state space: the gains
subspace is then the span of the assets' excess-return vectors, and the statement is the biconditional
\lean{hasEMM_multi_iff_not_hasArbitrage} (\file{Foundations/FTAPMultiState}). The multi-period theorem
is stated for one risky asset; with several assets only the generators of the gains subspace would
change, but that version is not in the library. The restriction to finitely many states is what keeps
the separating functional elementary and the measure a genuine probability vector; the next two
sections give it up.

%======================================================================
\section{One period, one asset: the elementary theorem}\label{sec:scalar}

Leaving finitely many states changes the problem. On an arbitrary probability space the attainable gains
need not form a closed set, and the separation argument of \cref{sec:finite} has nothing compact to
grip. For a single scalar return, though, the equivalent martingale measure can still be written down
directly, with no separation at all.

Fix a measurable $Y\colon\Omega\to\R$ on an arbitrary $(\Omega,P)$. Two difficulties replace the single
one of the finite case. First, $Y$ need not be integrable, so $\E_Q[Y]$ may fail even to make sense.
Second, a fair measure must be produced from $P$ rather than chosen from a simplex. Both yield to changes
of density.

Integrability is bought by \emph{tempering}. The bounded strictly positive weight $w=(1+\abs{Y})^{-1}$,
normalised by $\kappa=\E[w]\in(0,1]$, defines $\tilde P=w/\kappa\cdot P$, equivalent to $P$, under which
$Y$ is integrable because $\abs{Y}w/\kappa\le\kappa^{-1}$. No arbitrage is an almost-sure notion, so it
survives the passage to the equivalent $\tilde P$. Over $\tilde P$ the construction is a dichotomy on
the sign of $Y$: if $Y\ge0$ almost surely and is not almost surely zero then $\theta=1$ is an
arbitrage, and symmetrically for $Y\le0$, so under no arbitrage the positive and negative parts of $Y$
both carry mass. A two-region density that re-weights the two sides,
\[
  Z=\lambda\,\ind_{\{Y\ge0\}}+\mu\,\ind_{\{Y<0\}},\qquad \lambda,\mu>0,
\]
can then be balanced so that $\E[ZY]=0$ and normalised so that $\E[Z]=1$; one choice of $\lambda,\mu$,
explicit in the two tail integrals of $Y$, does both at once. If $Y=0$ almost surely, $Q=P$ will do. The
balancing core is \lean{exists_isEMM_of_pos_tails}; removing the integrability assumption by the tempering
above gives \lean{exists_isEMM_of_noArbitrage}, whose measure has $P$-density proportional to
$Z\,(1+\abs{Y})^{-1}$, with $Z$ computed under $\tilde P$. The forward direction closes the biconditional
(\cref{lst:scalar}).

\begin{leancode}[caption={The scalar one-period theorem (\texttt{Foundations/FTAPOnePeriod}); $Y:\Omega\to\mathbb{R}$ on an arbitrary probability space.},label=lst:scalar]
theorem ftap_one_period (hY : Measurable Y) :
    NoArbitrage P Y ↔ ∃ Q, IsEMM P Y Q
\end{leancode}

This is the one-asset case of the one-period theorem of Dalang, Morton and Willinger, in the form found
in Föllmer and Schied~\cite[Ch.~1]{FollmerSchied2016}. No Hahn--Banach theorem and no Kreps--Yan
argument enter; the measure is built rather than selected. What the argument leaves open is dimension.
A single asset has only two sides to balance, and the dichotomy that drives the construction does not
by itself reach a market of several assets.

%======================================================================
\section{One period, many assets: the variational construction}\label{sec:vector}

The $d$-asset one-period theorem is where the construction leaves the textbook proof furthest behind.
The return is now $Y\colon\Omega\to F$ with $F$ a finite-dimensional inner-product space, and a position
is a constant vector $\theta\in F$. The classical route separates the cone of attainable claims from the
non-negative payoffs in $L^1$ by a Hahn--Banach functional, an argument whose infinite-dimensional
duality and cone-closedness are heavy to formalize. We replace it: the equivalent martingale measure is
the minimiser of one smooth convex function on the finite-dimensional space $F$, and every step is a
finite-dimensional or elementary-integral fact. Throughout, the Lean spelling of the gain
$\langle\theta,Y\rangle$ is \lstinline|inner ℝ θ (Y ω)|.

\subsection{The potential}
For $\theta\in F$ set
\[
  f(\theta)=\E\bigl[\log\bigl(1+e^{\langle\theta,Y\rangle}\bigr)\bigr],
\]
the expected \emph{softplus} of the gain. The softplus $u\mapsto\log(1+e^u)$ is smooth and convex, with
derivative the logistic function $\sigma(u)=e^u/(1+e^u)\in(0,1)$, and it satisfies
\[
  u^+\;\le\;\log(1+e^u)\;\le\;\abs{u}+\log 2 .
\]
The upper bound dominates the integrand by $\norm{\theta}\norm{Y}+\log2$, so $f$ is finite and continuous
whenever $Y$ is integrable; the lower bound will give coercivity. Because $\sigma$ is bounded,
differentiation under the integral sign is dominated by $\norm{u}\norm{Y}$, and the derivative in a direction
$u$ is
\[
  \frac{\dd}{\dd t}\Big|_{t=0} f(\theta+tu)=\E\bigl[\sigma(\langle\theta,Y\rangle)\,\langle u,Y\rangle\bigr]
  \qquad(\text{\lean{hasDerivAt_potential_dir}}),
\]
the directional form of the gradient $\E[\sigma(\langle\theta,Y\rangle)\,Y]$, the $Y$-average weighted by the
logistic of the current gain (\cref{lst:potential}). The potential is also convex, but the proof never
uses that: existence of a minimiser needs only continuity and coercivity, and the first-order condition
holds at any minimiser.

\begin{leancode}[caption={The softplus potential (\texttt{Foundations/FTAPOnePeriodVector}).},label=lst:potential]
noncomputable def softplus (u : ℝ) : ℝ := Real.log (1 + Real.exp u)
noncomputable def logistic (u : ℝ) : ℝ := Real.exp u / (1 + Real.exp u)
noncomputable def potential (θ : F) : ℝ :=
  ∫ ω, softplus (inner ℝ θ (Y ω)) ∂P
\end{leancode}

\subsection{No arbitrage gives coercivity}
No arbitrage is what makes $f$ attain a minimum. Redundant directions are set aside first. The \emph{gains kernel}
\[
  N=\{\theta\in F:\ \langle\theta,Y\rangle=0\ \ P\text{-a.s.}\}
\]
collects the positions with almost surely zero gain. It is a linear subspace, and $f$ is constant along
it, since adding an element of $N$ changes no gain. On the orthogonal complement $N^\perp$ the picture is
rigid. The positive-gain average $g(\theta)=\E[\langle\theta,Y\rangle^+]$ is continuous and positively
homogeneous, and under no arbitrage it is strictly positive on $N^\perp\setminus\{0\}$: if $g(\theta)=0$
then $\langle-\theta,Y\rangle\ge0$ almost surely, so no arbitrage puts $\theta$ in $N$, while
$N\cap N^\perp=\{0\}$. Its minimum $c$ over the unit sphere of $N^\perp$ is therefore positive, and since
the softplus dominates the positive part,
\[
  f(\theta)\;\ge\;g(\theta)\;\ge\;c\,\norm{\theta}\qquad(\theta\in N^\perp).
\]
No arbitrage makes $f$ coercive on $N^\perp$ (\cref{lst:coercive}). The converse holds as well,
though the theorem does not need it and it is not formalized: if $\theta$ is an arbitrage, then at every
point $\eta$ the derivative of $f$ in the direction $-\theta$ is
$-\E[\sigma(\langle\eta,Y\rangle)\langle\theta,Y\rangle]<0$, so $f$ has no minimiser.

\begin{leancode}[caption={Coercivity from no arbitrage (\texttt{Foundations/FTAPOnePeriodVector}).},label=lst:coercive]
lemma exists_pos_lower_bound (hYint : Integrable Y P) (hNA : NoArbitrage P Y)
    (hNbot : (gainsKernel P Y)ᗮ ≠ ⊥) :
    ∃ c > 0, ∀ θ ∈ (gainsKernel P Y)ᗮ, c * ‖θ‖ ≤ potential P Y θ
\end{leancode}

\subsection{The minimiser and its measure}
Coercivity and continuity give a minimiser $\theta_0$ of $f$ on the closed subspace $N^\perp$, by
compactness of a closed ball (\lean{exists_global_min_potential}); because
$f$ is flat along $N$ and $F=N\oplus N^\perp$, the same $\theta_0$ minimises $f$ over all of $F$. At a
global minimiser every directional derivative vanishes, so
\[
  \E\bigl[\sigma(\langle\theta_0,Y\rangle)\,Y\bigr]=0
\]
(\cref{lst:foc}). The weight $z=\sigma(\langle\theta_0,Y\rangle)$ lies strictly between $0$ and $1$: it is
bounded, strictly positive, and makes $Y$ fair. The measure $Q$ with $P$-density $z/\E[z]$ is therefore a
probability measure equivalent to $P$, under which $Y$ is integrable with $\E_Q[Y]=0$. If $N^\perp$ is
trivial, every direction is redundant, $Y=0$ almost surely, and $Q=P$ will do.

\begin{leancode}[caption={The first-order condition (\texttt{Foundations/FTAPOnePeriodVector}).},label=lst:foc]
lemma integral_logistic_smul_eq_zero (hY : Measurable Y) (hYint : Integrable Y P)
    {θ₀ : F} (hmin : ∀ θ, potential P Y θ₀ ≤ potential P Y θ) :
    ∫ ω, logistic (inner ℝ θ₀ (Y ω)) • Y ω ∂P = 0
\end{leancode}

Two features set this apart from the textbook proof. No non-redundancy hypothesis is needed: redundant
assets live in $N$, the minimisation takes place on $N^\perp$, and the kernel is absorbed rather than
assumed away. And integrability of $Y$ is not assumed: as in \cref{sec:scalar}, the bounded tempering
$(1+\norm{Y})^{-1}$ reduces to the integrable case, and equivalence of measures carries the result back.
The minimisation then takes place under the tempered measure $\tilde P$, so in general the density
of $Q$ with respect to $P$ is proportional to $\sigma(\langle\theta_0,Y\rangle)\,(1+\norm{Y})^{-1}$, with $\theta_0$
a minimiser of the potential under $\tilde P$.
The backward direction is \lean{exists_isEMM_of_noArbitrage}; with the forward direction it is the
biconditional (\cref{lst:vector}).

\begin{leancode}[caption={The $d$-asset one-period theorem (\texttt{Foundations/FTAPOnePeriodVector}); the context is that of \cref{lst:defs}.},label=lst:vector]
theorem ftap_one_period_vector (hY : Measurable Y) :
    NoArbitrage P Y ↔ ∃ Q, IsEMM P Y Q
\end{leancode}

Where the classical proof invokes the dual of $L^\infty$ and the closedness of a cone in $L^0$, this one
invokes coercivity, a minimum over a sphere, and differentiation under an integral, each already in
Mathlib. The reward is an equivalent martingale measure given by a formula in $\theta_0$,
$\sigma(\langle\theta_0,Y\rangle)$ up to normalisation and tempering, rather than produced by a separating
functional. The minimiser itself is obtained by compactness, not computed; ``explicit'' here means a
formula in $\theta_0$, not a closed form. The formula lives in the proof: the theorem asserts existence,
and what the library states about the formula is the first-order condition of \cref{lst:foc}, from which
the martingale property of $Q$ follows by normalisation.

\subsection{Why the softplus}\label{sec:why-softplus}
Variational proofs of the one-period theorem are not new: Rogers~\cite{Rogers1994} obtains the
equivalent martingale measure from a minimisation argument rather than a separation. What matters for a
formalization is the choice of potential, and the softplus is chosen for what the argument asks of it.
For a convex $C^1$ function $\phi$ in place of the softplus, the proof uses four properties:
\begin{enumerate}[label=(\roman*)]
\item $\phi(u)\ge u^+$, which turns no arbitrage into coercivity;
\item $\abs{\phi(u)}\le a\abs{u}+b$, which makes $\E[\phi(\langle\theta,Y\rangle)]$ finite for integrable $Y$;
\item $\phi'$ bounded, which dominates the derivative under the integral by an $L^1$ function;
\item $\phi'>0$, which makes the weight $\phi'(\langle\theta_0,Y\rangle)$ a strictly positive density.
\end{enumerate}
The softplus has all four, with $\phi'=\sigma\in(0,1)$. The exponential $\phi(u)=e^u$ has (i) and (iv)
but neither (ii) nor (iii): its minimiser gives the Esscher density proportional to
$e^{\langle\theta_0,Y\rangle}$, which needs exponential moments of $Y$, and the bounded tempering
$(1+\norm{Y})^{-1}$ supplies only integrability. The logistic density is thus the softplus analogue of
the Esscher measure, not the Esscher measure itself (the first version of this paper conflated the
two). It is bounded, and it exists whenever $Y$ is integrable.

\begin{remark}[interpretation, not formalized]
The convex conjugate of the softplus is $q\log q+(1-q)\log(1-q)$ on $[0,1]$, minus the binary entropy
$H(q)$. Interchanging the minimum over $\theta$ with the supremum over weights, which convex duality
justifies but we do not formalize, identifies $\min_\theta f(\theta)$ with
$\max\{\E[H(z)]:\ 0\le z\le1,\ \E[zY]=0\}$, attained at $z=\sigma(\langle\theta_0,Y\rangle)$. The density
the construction produces is, up to normalisation, the fair weight of maximal expected binary entropy.
\end{remark}

%======================================================================
\section{A worked example}\label{sec:example}

A two-state market makes the construction concrete and shows the gains kernel at work. The example is
computed by hand; it illustrates the Lean argument and is not itself a Lean theorem. Let
$\Omega=\{+,-\}$ with $P(+)=P(-)=\tfrac12$, and take two assets whose discounted excess returns are
\[
  Y_1(+)=a,\quad Y_1(-)=-b\quad(a,b>0),\qquad Y_2=c\,Y_1\quad(c\neq0),
\]
so the second asset is redundant. A position $\theta=(\theta_1,\theta_2)$ has gain
$\langle\theta,Y\rangle=(\theta_1+c\,\theta_2)\,Y_1$, which depends on $\theta$ only through the scalar
$s=\theta_1+c\,\theta_2$. The gains kernel is the line $N=\{\theta:\theta_1+c\,\theta_2=0\}$, and the
potential collapses to a function of $s$ alone,
\[
  f(\theta)=\tfrac12\log(1+e^{sa})+\tfrac12\log(1+e^{-sb}).
\]
Its derivative in $s$ is $\tfrac12\bigl(a\,\sigma(sa)-b\,\sigma(-sb)\bigr)$, strictly increasing in $s$,
negative for $s\to-\infty$ and positive for $s\to+\infty$, so it vanishes at a unique $s_0$. The logistic
weight $z=\sigma(s_0Y_1)$, normalised, is the density, and the first-order condition
$a\,\sigma(s_0a)=b\,\sigma(-s_0b)$ is $\E_Q[Y_1]=0$. In two states that pins the measure to
\[
  Q(+)=\frac{b}{a+b},\qquad Q(-)=\frac{a}{a+b},
\]
the classical risk-neutral probability, which is also the explicit weight of the two-state module
(\lean{emmWeightUp}, after~\cite{Nagy2026}). The redundant asset adds no constraint: once $Y_1$ is fair,
$Y_2=c\,Y_1$ is fair automatically, as the minimisation over $N^\perp$ predicts. In higher dimension the
first-order condition is no longer a scalar equation, but the structure is the same: the kernel absorbs
the redundant directions, and for integrable returns the logistic of the minimising gain is the density.

%======================================================================
\section{One separation for pricing and risk}\label{sec:separation}

The root of \cref{lst:cone} has a sibling. A single point $x_0$ outside a closed convex cone $C$ that
contains the origin is separated from it by a functional $p$ with $\sum_ip_iv_i\le0$ on $C$ and
$\sum_ip_ix_{0,i}>0$ (\lean{exists_separating_of_not_mem_cone}). Both are derived in the same file from
Mathlib's geometric Hahn--Banach theorems through two shared lemmas, a basis expansion of a functional and
the fact that a functional bounded below on a cone is non-negative on it. The first is the pricing
reading of separation; the second is the risk reading.

\paragraph{Risk}
A risk measure $\rho$ on a finite state space is \emph{coherent} if it is monotone, cash-invariant,
positively homogeneous and subadditive~\cite{ADEH1999}. Its acceptance set $\{X:\rho(X)\le0\}$ is then a
closed convex cone, and a position that is not acceptable is a point outside it. Separating the rejected
position $X+(\rho(X)-\varepsilon)\ind$ from the acceptance cone and normalising the separating functional
yields a probability vector that prices every acceptable position non-negatively; letting
$\varepsilon\downarrow0$ gives the representation of Artzner, Delbaen, Eber and Heath
(\cref{lst:risk}): $\rho(X)$ is the least upper bound of the expected losses $\E_q[-X]$ over the
representing set.

\begin{leancode}[caption={The coherent-risk representation on a finite state space (\texttt{RiskMeasures/AcceptanceSet}).},label=lst:risk]
def representingSet {ι : Type*} [Fintype ι] (ρ : (ι → ℝ) → ℝ) :
    Set (ι → ℝ) :=
  {q | (∀ i, 0 ≤ q i) ∧ (∑ i, q i = 1) ∧
    ∀ Z ∈ {Y : ι → ℝ | ρ Y ≤ 0}, 0 ≤ ∑ i, q i * Z i}

theorem coherentRisk_isLUB {ι : Type*} [Fintype ι] {ρ : (ι → ℝ) → ℝ}
    (hρ : IsCoherentRisk ρ) (X : ι → ℝ) :
    IsLUB ((fun q ↦ ∑ i, q i * (- X i)) '' representingSet ρ) (ρ X)
\end{leancode}

So the equivalent martingale measure of a finite market and the representing scenarios of a coherent risk
measure are the same kind of object, produced the same way: a strictly positive (respectively
non-negative) normalised functional that separates a cone. In the library this is more than an analogy.
Both theorems import the one module, and the two separations they rest on, cone from simplex and point
from cone, are proved from the same two helper lemmas; they are siblings, not one lemma applied twice. The statement of \lean{coherentRisk_isLUB} is also
restated with Mathlib alone and checked against the library by an external tool (\cref{sec:artifact}).

\paragraph{Superhedging}
Superhedging shows where the separation has not yet been carried. Weak duality needs none: every
martingale measure, equivalent or not, prices a claim below the cost of any portfolio that dominates it
(\lean{emm_le_superReplication}), by a direct calculation. Strong duality, that the superhedging price
equals the supremum of martingale prices, is the separation-based converse. It needs the cone of
super-replicable payoffs to be closed, which for a finitely generated cone is classical but is not
available in Mathlib at the library's pin; Mathlib's cone-separation theorem itself is. Strong duality is
not formalized.

%======================================================================
\section{What is proved, and what is not}\label{sec:boundary}

The theorems above are complete and machine-checked: the finite-state one-period theorem with several
assets, the finite-state multi-period theorem with one asset, and the one-period theorem on an arbitrary
probability space with one asset and with several. Each is a biconditional between no arbitrage and the
existence of an equivalent martingale measure (\cref{fig:ladder}).

\begin{figure}[tbp]
\centering
\begin{tikzpicture}[
  node distance=7mm and 6mm,
  box/.style={draw=mfaccent!70, fill=mfbg, rounded corners=2pt, align=center,
    text width=64mm, inner sep=4pt, font=\small},
  open/.style={box, dashed, fill=white},
  arr/.style={-{Stealth[length=2.2mm]}, draw=mfaccent},
  lab/.style={font=\footnotesize\itshape, text=mfmuted}]
\node[box] (a) {Finite states, one period, $d$ assets\\[1pt]
  \texttt{\footnotesize hasEMM\_multi\_iff\_not\_hasArbitrage}};
\node[box, below=of a] (b) {Finite states, $T$ periods, one asset\\[1pt]
  Harrison--Pliska $\cdot$ \texttt{\footnotesize ftap\_discrete}};
\node[box, below=of b] (c) {Arbitrary $\Omega$, one period, one asset\\[1pt]
  \texttt{\footnotesize ftap\_one\_period}};
\node[box, below=of c] (d) {Arbitrary $\Omega$, one period, $d$ assets\\[1pt]
  variational EMM $\cdot$ \texttt{\footnotesize ftap\_one\_period\_vector}};
\node[open, below=of d] (e) {Arbitrary $\Omega$, $T$ periods\\[1pt]
  Dalang--Morton--Willinger $\cdot$ forward direction only};
\node[open, below=of e] (f) {Continuous time\\[1pt]
  Delbaen--Schachermayer $\cdot$ forward direction only};
\draw[arr] (a) -- node[right, lab] {add periods} (b);
\draw[arr] (b) -- node[right, lab] {drop finiteness} (c);
\draw[arr] (c) -- node[right, lab] {add assets} (d);
\draw[arr, dashed] (d) -- node[right, lab] {add dynamics} (e);
\draw[arr, dashed] (e) -- node[right, lab] {continuous trading} (f);
\end{tikzpicture}
\caption{The settings formalized here (solid) and those that remain open or partial (dashed). Each solid
box is a machine-checked biconditional between no arbitrage and an equivalent martingale measure.}
\label{fig:ladder}
\end{figure}

\paragraph{The general multi-period theorem}
One rung of the ladder is missing its converse. The forward direction holds on an arbitrary probability
space over several periods: an equivalent martingale measure excludes arbitrage by bounded predictable
strategies in one risky asset (\lean{emm_implies_no_arbitrage}). The theorem of Dalang, Morton and Willinger places a
multi-period market on an arbitrary probability space: trading is dynamic, strategies are predictable
processes, and the attainable gains live in $L^0$ rather than in a finite-dimensional space. Its classical
proofs need two ingredients that the cases here sidestep: the closedness of the cone of attainable claims
in $L^0$ under no arbitrage, and a measurable selection that assembles the one-period conditional
measures into a single density across time~\cite{DMW1990,KabanovStricker2001}. Both are substantial, and
neither is formalized; the backward direction is open. So is the finite-state multi-period theorem with
several assets. The finite-state case captures the multi-period structure when the sample space is
finite; the one-period cases capture the measure-theoretic difficulty in a single step. A natural route
to the general theorem through the present construction would apply the softplus minimisation
conditionally on $\F_{t}$, one period at a time; making the conditional minimiser measurable in $\omega$
is where a measurable-selection argument re-enters, and it is not attempted here.

\paragraph{Continuous time}
In continuous time the library proves the forward direction for simple strategies: if $Q\sim P$ and the
discounted price is a $Q$-martingale, then no strategy with finitely many trading dates and bounded
predictable holdings is an arbitrage (\lean{isEMM_noArbitrageSimple}, \file{Foundations/ContinuousMarket}).
For the Black--Scholes model the equivalent martingale measure is not assumed but constructed, by a
Girsanov change of measure~\cite{CoelhoLibrary,CoelhoIto}. The converse, that no free lunch with
vanishing risk yields an equivalent ($\sigma$-)martingale measure~\cite{DelbaenSchachermayer1994}, is not
formalized.

\paragraph{The second theorem}
Around these results the library carries the uniqueness half of first-course arbitrage theory. In the
Brownian market with a driftless price $S=S_0+\int\sigma\,\dd B$, $\sigma\neq0$ almost everywhere, a
probability measure with square-integrable density under which $S$ is a martingale agrees with $P$ on
$\F^B_T$ (\lean{measure_eq_of_density}, \file{Foundations/PricingMeasureL2Density}). Its proof runs
through martingale representation and belongs to the It\^o calculus~\cite{CoelhoIto}. It is the direction
from completeness to uniqueness; the converse, which needs the extreme-point characterisation of Jacod and
Yor, is not formalized.

%======================================================================
\section{The artifact}\label{sec:artifact}

The development is part of a Lean~4 library that builds against pinned versions of its dependencies:
the Lean toolchain \texttt{v4.33.0-rc1}, Mathlib at revision \texttt{0434c03}, and the
\texttt{BrownianMotion} package~\cite{BrownianMotionRepo} at \texttt{0d5b6eb}, which the results of this
paper do not use. The state described here is release \texttt{v1.5.0} (commit \texttt{c0180c8}),
archived on Zenodo~\cite{Zenodo}. A plain \texttt{lake build} from a clean checkout is the canonical
check.

Honesty about what the kernel accepts is enforced rather than asserted. Two files run
\texttt{\#print axioms} under \texttt{\#guard\_msgs}: a curated audit of headline results, and a generated one
that pins every library constant a benchmark proof cites, resolved by declaration rather than by name.
Each theorem discussed here depends only on \texttt{propext}, \texttt{Classical.choice} and
\texttt{Quot.sound}; a \texttt{sorry} or an additional axiom anywhere beneath one of them fails the build.
A verification ledger records, for every benchmark entry, a hash of the exact sources and toolchain pins it
was last checked against, so that an entry whose inputs change cannot keep a stale verdict. For the
coherent-risk theorem of \cref{sec:separation} the repository also carries a Mathlib-only restatement,
paired with a witness that its hypothesis is satisfiable, and a proof of both from the library that every
build re-checks; the external checker Comparator~\cite{Comparator} is configured to confirm that the proved
statements are exactly the restated ones and that only the three standard axioms are used. The results discussed are collected in
\cref{tab:results}.

\begin{table}[tbp]
\centering\footnotesize
\begin{tabularx}{\linewidth}{@{}>{\raggedright\arraybackslash}p{0.40\linewidth}>{\raggedright\arraybackslash}X>{\raggedright\arraybackslash}p{0.26\linewidth}@{}}
\toprule
Lean declaration & Module & Benchmark entry \\
\midrule
\lean{MathFin.ftap_discrete} & \file{Foundations/FTAPDiscrete} & \bench{mf-ftap-discrete-complete} \\
\lean{MathFin.hasEMM_multi_iff_not_hasArbitrage} & \file{Foundations/FTAPMultiState} & \bench{mf-ftap-single-period-complete} \\
\lean{MathFin.OnePeriod.ftap_one_period} & \file{Foundations/FTAPOnePeriod} & \bench{mf-ftap-one-period-general} \\
\lean{MathFin.OnePeriodVector.ftap_one_period_vector} & \file{Foundations/FTAPOnePeriodVector} & \bench{mf-ftap-one-period-vector} \\
\lean{MathFin.exists_pos_separating_of_cone_disjoint_simplex} & \file{Foundations/ConvexDuality} & \bench{mf-convex-duality-root} \\
\lean{MathFin.coherentRisk_isLUB} & \file{RiskMeasures/AcceptanceSet} & \bench{mf-coherent-risk-representation} \\
\lean{MathFin.emm_le_superReplication} & \file{Foundations/SuperhedgingDuality} & \bench{mf-superhedging-emm-bound} \\
\lean{MathFin.emm_implies_no_arbitrage} & \file{Foundations/FTAP} & \bench{mart-thm-2.6.7} \\
\lean{MathFin.ContinuousMarket.isEMM_noArbitrageSimple} & \file{Foundations/ContinuousMarket} & \bench{gir-continuous-emm-forward} \\
\lean{MathFin.PricingMeasureL2Density.measure_eq_of_density} & \file{Foundations/PricingMeasureL2Density} & \bench{gir-pricing-measure-density} \\
\bottomrule
\end{tabularx}
\caption{The results discussed in this paper, their modules under \texttt{MathFin/}, and their benchmark
entries. Every declaration is pinned by the axiom audits to \texttt{propext}, \texttt{Classical.choice} and
\texttt{Quot.sound}.}
\label{tab:results}
\end{table}

Within the wider library these sit among 535 catalogued results: 504 complete formalizations, 18 thin
wrappers over Mathlib or \texttt{BrownianMotion} theorems, and 13 reduced cores that prove less than the
result they are named after, with the 522 complete-or-wrapper entries counting as delivered. The catalogue
is published as a dataset in which every entry carries its statement, its status and a description of
what its formalization leaves out~\cite{HFDataset}. The library is released under the Apache~2.0
licence.

\paragraph{Provenance}
The Lean development was written by the author with an AI coding assistant (Claude Code), as recorded in
the repository's provenance file; every proof is checked by the Lean kernel, and the gates above apply
to it unchanged. This manuscript was revised with the same assistant, and its statements were checked
against the Lean sources of release \texttt{v1.5.0}.

%======================================================================
\begin{thebibliography}{99}
\small
\bibitem{ADEH1999} P.~Artzner, F.~Delbaen, J.-M.~Eber, D.~Heath. Coherent measures of risk.
  \emph{Mathematical Finance}, 9(3):203--228, 1999.
\bibitem{AvigadHolzlSerafin2017} J.~Avigad, J.~H\"olzl, L.~Serafin. A formally verified proof of the
  central limit theorem. \emph{Journal of Automated Reasoning}, 59(4):389--423, 2017.
\bibitem{BrownianMotionRepo} R.~Degenne et al. \emph{brownian-motion}: a Lean~4 development of Brownian
  motion and stochastic integrals. \url{https://github.com/RemyDegenne/brownian-motion}, 2024--2026.
\bibitem{CoelhoIto} R.~Coelho. A machine-checked It\^o calculus for Brownian motion.
  arXiv:2606.15089, 2026.
\bibitem{CoelhoLibrary} R.~Coelho. A formally verified library of mathematical finance in Lean~4.
  arXiv:2606.01356, 2026.
\bibitem{Comparator} \emph{Comparator}: a checker for Lean proofs against Mathlib-only statements.
  \url{https://github.com/leanprover/comparator}, 2026.
\bibitem{DMW1990} R.~C.~Dalang, A.~Morton, W.~Willinger. Equivalent martingale measures and no-arbitrage
  in stochastic securities market models. \emph{Stochastics and Stochastic Reports}, 29(2):185--201,
  1990.
\bibitem{DegenneEtAl2025} R.~Degenne, D.~Ledvinka, E.~Marion, P.~Pfaffelhuber. Formalization of Brownian
  motion in Lean. arXiv:2511.20118, 2025.
\bibitem{DelbaenSchachermayer1994} F.~Delbaen, W.~Schachermayer. A general version of the fundamental
  theorem of asset pricing. \emph{Mathematische Annalen}, 300(1):463--520, 1994.
\bibitem{deMouraUllrich2021} L.~de~Moura, S.~Ullrich. The Lean~4 theorem prover and programming
  language. In \emph{CADE~28}, LNCS 12699, pp.~625--635, 2021.
\bibitem{EchenimGuiolPeltier2020} M.~Echenim, H.~Guiol, N.~Peltier. Formalizing the Cox--Ross--Rubinstein
  pricing of European derivatives in Isabelle/HOL. \emph{Journal of Automated Reasoning}, 64:737--765,
  2020.
\bibitem{FollmerSchied2016} H.~F\"ollmer, A.~Schied. \emph{Stochastic Finance: An Introduction in
  Discrete Time}, 4th ed. De Gruyter, 2016.
\bibitem{HarrisonKreps1979} J.~M.~Harrison, D.~M.~Kreps. Martingales and arbitrage in multiperiod
  securities markets. \emph{Journal of Economic Theory}, 20(3):381--408, 1979.
\bibitem{HarrisonPliska1981} J.~M.~Harrison, S.~R.~Pliska. Martingales and stochastic integrals in the
  theory of continuous trading. \emph{Stochastic Processes and their Applications}, 11(3):215--260,
  1981.
\bibitem{HFDataset} R.~Coelho. \emph{formal-mathfin-theorems} (dataset).
  \url{https://huggingface.co/datasets/formal-applied-math/formal-mathfin-theorems}, 2026.
\bibitem{KabanovStricker2001} Y.~Kabanov, C.~Stricker. A teachers' note on no-arbitrage criteria. In
  \emph{S\'eminaire de Probabilit\'es XXXV}, LNM 1755, pp.~149--152. Springer, 2001.
\bibitem{Keskin2023} A.~Keskin. A formalization of martingales in Isabelle/HOL. arXiv:2311.06188, 2023.
\bibitem{Mathlib2020} The mathlib Community. The Lean mathematical library. In \emph{CPP 2020},
  pp.~367--381, 2020.
\bibitem{Nagy2026} T.~Nagy. From It\^o to Black--Scholes: a machine-verified derivation in Lean~4. SSRN
  Working Paper 6336503, 2026.
\bibitem{Rogers1994} L.~C.~G.~Rogers. Equivalent martingale measures and no-arbitrage.
  \emph{Stochastics and Stochastic Reports}, 51(1--2):41--49, 1994.
\bibitem{Zenodo} R.~Coelho. \emph{formal-mathfin: a formally verified library of mathematical finance in
  Lean~4} (software). Zenodo, \url{https://doi.org/10.5281/zenodo.20477781}, 2026.
\end{thebibliography}

\end{document}
